\section{Note on sources and delivery}

For an adult parish assembly on the First Sunday of Advent, Year B, 29 November 2026. The 1,121 spoken words take about 10 minutes 40 seconds at 105 words per minute. A silent reading checked the change of scene from the Mount of Olives to Gethsemane and the invitation to keep someone company. The timing is an estimate.

The homily follows the studies' interpretation \textit{The people who ask God to turn them}. Mark records the warning to everyone and then the request made to three particular friends. Jesus' instruction to pray against temptation occurs in the garden (14:38). Jerome explains why Isaiah's speakers can still call on their Father. Augustine explains why the people asking to invoke God ask him first for life. The connection with the Collect, Communion and a promised visit is the homily's application.

\section{References}
\begin{enumerate}
\item U.S. \textit{Lectionary for Mass}, vol.~I, no.~2. Mark 13:33--37, Isaiah 63:16b--17, 19b and 64:2--7, and Psalm 80 (79):2--3, 15--16, 18--19 with response from v.~4. See the \href{https://bible.usccb.org/bible/readings/112926.cfm}{U.S. bishops' readings for 29 November 2026}. Narrative context in Mark 13:1--4 and 14:29--42. Israel's vine brought from Egypt in Psalm 80 (79):9.
\item Jerome, \href{https://www.documentacatholicaomnia.eu/03d/0347-0420,_Hieronymus,_In_Isaiam_Pars_4,_LT.pdf}{\textit{Commentary on Isaiah}}, book XVII, Latin part IV, pp.~158--159, on Isaiah 64:7--8 in the Vulgate.
\item Augustine of Hippo, \href{https://www.newadvent.org/fathers/1801080.htm}{\textit{Exposition on Psalm} 79}, §§3, 11, English heading Psalm 80, \textit{Nicene and Post-Nicene Fathers}, first series, vol.~8.
\item \textit{Roman Missal}, Third Edition for the United States (2011), First Sunday of Advent, Collect, p.~139. United States Conference of Catholic Bishops, \href{https://www.usccb.org/prayer-and-worship/liturgical-year-and-calendar/advent/commentary-on-advent-proper-prayers}{\textit{Commentary on the Proper Prayers of Advent from the Roman Missal}}, First Sunday, Collect, on meeting Christ at Communion and in people who need our service.
\end{enumerate}
